\section{Do the Assumptions Hold? Calibration and Tests}
\label{sec:assumptions}

\subsection{Measuring Fill Intensity from Bars}

Bars do not show the order book, so we infer fills from price excursions. Suppose that at the end
of minute $t$ a dealer posts an ask at $S_t+\delta$, where $S_t$ is the close. The ask is
\emph{filled} if the next minute's high reaches it, $H_{t+1}\ge S_t+\delta$; the bid at
$S_t-\delta$ is filled if $L_{t+1}\le S_t-\delta$. For each distance $\delta$ on a grid we count the
fraction $p(\delta)$ of minutes with a fill. If fills arrive as a Poisson process with intensity
$\lambda(\delta)$, the probability of at least one fill in one minute is $1-e^{-\lambda(\delta)}$, so
\begin{equation}
\lambda(\delta) = -\ln\big(1-p(\delta)\big), \qquad
\ln\lambda(\delta) = \ln A - k\,\delta .
\label{eq:fit}
\end{equation}
Under A1 the points $(\delta,\ln\lambda)$ lie on a straight line, and an ordinary least-squares fit
gives $\ln A$ as the intercept and $-k$ as the slope. We fit only distances with
$0.02<p(\delta)<0.98$, where $p$ is estimated precisely. The volatility $\sigma$ is the standard
deviation of one-minute changes in the close. All three are estimated on the calibration window
only. A synthetic test in the code checks that this procedure recovers known values of $A$ and $k$
to within 3\%.

\subsection{A1--A2: Fill Intensity Is Exponential and Symmetric}

\begin{figure}[!htbp]
\centering
\includegraphics[width=\textwidth]{calib_fill_curves}
\caption{The model's central assumption, tested. Implied fill intensity $\lambda(\delta)$ against
quote distance on a log scale, calibration window. Filled circles: ask side; open squares: bid side;
grey points lie outside the fitting band. A straight line on this scale is exactly
$\lambda=Ae^{-k\delta}$.}
\label{fig:fill}
\end{figure}

Figure~\ref{fig:fill} is the central test of the paper, and the answer is clear. On a log scale
the points lie on a straight line for both stocks and both sides of the book, with
$R^2\ge\RELIANCERTwo$ for RELIANCE and $R^2\ge\KSCLRTwo$ for KSCL (Table~\ref{tab:calibration}).
The two sides give nearly the same slope ($k_{\text{bid}}/k_{\text{ask}}$ within 10\%), so one
curve for both sides (A2) is a fair simplification. The liquid stock has the higher $A$ (more fills
at any distance) and the higher $k$ (its book thins out faster, because competing quotes sit close
to the mid).

The residuals in Figure~\ref{fig:residuals} show where the fit is imperfect. They curve gently
upward at both ends: quotes at the mid are filled a little more often than the line predicts,
and so are quotes far away. The tail of real fill intensity is therefore slightly fatter than
exponential, consistent with occasional large orders sweeping several price levels. The
deviations are small: apart from the single point at $\delta=0$, they stay within about 0.25 in log units, so A1 holds.

\begin{figure}[!htbp]
\centering
\includegraphics[width=\textwidth]{calib_fill_residuals}
\caption{Residuals of the log-linear fit (circles: ask, squares: bid). The gentle upward curve at
both ends means real fills have a slightly fatter tail than $e^{-k\delta}$.}
\label{fig:residuals}
\end{figure}

\begin{table}[!htbp]
\centering
\caption{Calibrated parameters, first \CalibDays{} days. The model uses one curve for both sides,
so $A$ and $k$ are averaged over the bid and the ask.}
\label{tab:calibration}
\small
\begin{tabular}{lrr}
\toprule
Parameter & RELIANCE & KSCL \\
\midrule
$\sigma$ (INR$/\sqrt{\text{min}}$) & 0.676 & 0.863 \\
$A$ (fills per minute), ask / bid & 1.341 / 1.631 & 0.756 / 0.927 \\
$k$ (per INR), ask / bid & 2.400 / 2.636 & 1.747 / 1.863 \\
$A$ used (mean of sides) & 1.486 & 0.842 \\
$k$ used (mean of sides) & 2.518 & 1.805 \\
$R^2$ of log-linear fit, ask / bid & 0.984 / 0.989 & 0.992 / 0.995 \\
Points in fit, ask / bid & 51 / 48 & 46 / 46 \\
\bottomrule
\end{tabular}

\end{table}

\subsection{A3--A4: The Mid-Price Is Close to a Random Walk, Except for the Bounce}

The mean one-minute price change is statistically zero for both stocks, so A3 holds. A4 bends.
Figure~\ref{fig:returns} shows that one-minute changes are far from normal: excess kurtosis is
\RELIANCEKurt{} for RELIANCE and \KSCLKurt{} for KSCL, because most minutes move by zero or one tick
and a few move by many (Figure~\ref{fig:qq}). Fat tails matter little for the model, whose
formulas only use the variance $\sigma^2$.

\begin{figure}[!htbp]
\centering
\includegraphics[width=\textwidth]{data_return_distribution}
\caption{Distribution of one-minute price changes, binned in whole ticks, on a log scale, with the
normal density of the same mean and variance. The empirical tails are far heavier.}
\label{fig:returns}
\end{figure}

\begin{figure}[!htbp]
\centering
\includegraphics[width=0.56\textwidth]{data_qq}
\caption{Normal quantile--quantile plot of standardised one-minute price changes. Both stocks
leave the line beyond about two standard deviations.}
\label{fig:qq}
\end{figure}

The important departure is in the dependence of price changes. For KSCL the first-order
autocorrelation is \KSCLAcfOne{} (Figure~\ref{fig:acf}), well outside the 95\% band, while
RELIANCE's is \RELIANCEAcfOne{}. This is the bid--ask bounce of Roll~\cite{roll1984}: when the last
trade is at the ask, the next is more likely at the bid, so recorded prices zig-zag even if the
true value does not move. The variance ratio of Lo and MacKinlay~\cite{lo1988},
$\mathrm{VR}(h)=\mathrm{Var}(S_{t+h}-S_t)/\big(h\,\mathrm{Var}(S_{t+1}-S_t)\big)$, which equals one for
a random walk, tells the same story over longer horizons (Figure~\ref{fig:vr}): it falls to
\KSCLVRfifteen{} at fifteen minutes for KSCL but stays at \RELIANCEVRfifteen{} for RELIANCE. The
practical consequence is that one-minute $\sigma$ overstates the true price risk of holding KSCL
over longer horizons. The model therefore over-penalises KSCL inventory, which makes it, if anything,
more conservative.

\begin{figure}[!htbp]
\centering
\includegraphics[width=\textwidth]{diag_autocorrelation}
\caption{Autocorrelation of one-minute price changes, lags 1--10 (shaded: 95\% band under
independence). KSCL's large negative first lag is the bid--ask bounce.}
\label{fig:acf}
\end{figure}

\begin{figure}[!htbp]
\centering
\includegraphics[width=\textwidth]{diag_variance_ratio}
\caption{Variance ratio against horizon. A random walk stays on the dashed line; values below one
mean price changes partly reverse. The illiquid stock reverses much more.}
\label{fig:vr}
\end{figure}

\subsection{A5: Parameters Are Not Constant}

Volatility is the assumption that breaks. Figure~\ref{fig:sigma} shows the fifteen-minute
volatility profile: $\sigma$ in the first fifteen minutes is two to three times its midday value, and across
the day the highest bucket is \RELIANCESigmaRatio{} times the lowest for RELIANCE (\KSCLSigmaRatio{}
for KSCL). The model uses one $\sigma$ for the whole day, so it under-prices inventory risk at the
open and over-prices it at midday. Across days (Figure~\ref{fig:stability}) the parameters are
steadier but not fixed: the daily $k$ varies with a coefficient of variation of about 0.17, and the
test window is somewhat more volatile than the calibration window, a realistic handicap for an
out-of-sample test.

\begin{figure}[!htbp]
\centering
\includegraphics[width=\textwidth]{data_intraday_sigma}
\caption{(a) Volatility in each fifteen-minute bucket relative to the session average; a model with
constant $\sigma$ would be the dashed line. (b) Share of minutes with at least one trade.}
\label{fig:sigma}
\end{figure}

\begin{figure}[!htbp]
\centering
\includegraphics[width=\textwidth]{calib_stability}
\caption{Parameters re-estimated day by day. Dashed lines: the calibration-window values used in
the rest of the paper. Shaded: calibration window.}
\label{fig:stability}
\end{figure}

\subsection{Scorecard}

Table~\ref{tab:scorecard} collects the verdicts. The assumption the model is built on, the
exponential fill curve, holds. The description of the price process bends, mostly for the illiquid
stock. Constant intraday volatility breaks. None of these failures is large enough to make the
calibrated model meaningless, so we now take it to the data.

\begin{table}[!htbp]
\centering
\caption{Assumption scorecard. ``Holds'': the data match the assumption. ``Bends'': a measurable
departure that the model tolerates. ``Breaks'': a departure large enough to change the quotes.}
\label{tab:scorecard}
\small
\begin{tabular}{llrrl}
\toprule
Assumption & Statistic & RELIANCE & KSCL & Verdict \\
\midrule
Fill intensity is $Ae^{-k\delta}$ & min $R^2$ & 0.984 & 0.992 & Holds \\
One curve for both sides & $k_{\text{bid}}/k_{\text{ask}}$ & 1.10 & 1.07 & Holds \\
Mid has no drift & $t$-stat of mean change & -0.26 & -0.43 & Holds \\
Independent increments & lag-1 autocorrelation & -0.014 & -0.123 & Bends (illiquid) \\
Random-walk scaling & VR(15) & 0.90 & 0.65 & Bends (illiquid) \\
Gaussian increments & excess kurtosis & 10.9 & 99.4 & Bends \\
Constant $\sigma$ through the day & max/min 15-min $\sigma$ & 3.1 & 3.2 & Breaks \\
Constant $k$ across days & coef.\ of variation of daily $k$ & 0.17 & 0.18 & Bends \\
\bottomrule
\end{tabular}

\end{table}
